\chapter{Энергосохраняющее четырёхтактное переключение и запрет невырожденного холодильника}

Предыдущая проверка доказала существование полной унитарной орбиты, на которой
семейный триплет охлаждается за счёт четырёхсостояний носителя. Но
произвольный унитарий может скрывать внешнюю работу. Теперь вводится самое
слабое физическое требование автономности:
\begin{equation}
 [W,H_S+H_C]=0.
 \label{eq:v6-energy-preserving-clock-control}
\end{equation}
Проверяется минимальная резонансная модель, сохраняющая точную одноосную
$\mathbb{RP}^2$-симметрию триплета и четыре различные фазы часов.

\section{Минимальные гамильтонианы}

Одноосное состояние с двукратно вырожденным поперечным уровнем является
гиббсовским для
\begin{equation}
 H_S=\varepsilon\operatorname{diag}(0,1,1).
 \label{eq:v6-uniaxial-qutrit-hamiltonian}
\end{equation}
В точке сосуществования требуемый разрыв равен
\begin{equation}
 \varepsilon_*
 =\frac1{\beta_c}\log\frac{2a_*}{1-a_*}
 =1.9664146236\ldots.
 \label{eq:v6-coexistence-gibbs-gap}
\end{equation}

Канонический порядок четырёх фаз задаёт невырожденную лестницу
\begin{equation}
 H_C=\varepsilon_*\operatorname{diag}(0,1,2,3),
 \label{eq:v6-four-tick-energy-ladder}
\end{equation}
после выбора длительности одного такта
\begin{equation}
 \tau_* =\frac{\pi}{2\varepsilon_*}
 =0.7988123705\ldots.
\end{equation}
Эта подстройка ещё не выведена из родителя, но даёт максимально
благоприятный резонансный тест.

\section{Блочное ограничение сохранения энергии}

Пусть часы начинаются в произвольном чистом состоянии. Обозначим через
$p_n$ вероятность их энергии $n\varepsilon_*$, $n=0,1,2,3$. Когерентности
между разными полными энергиями не влияют на населённость основного уровня
системы, поскольку его проектор коммутирует с $H_S+H_C$.

При полной энергии $n\varepsilon_*$ для $n=1,2,3$ существуют один вектор
\begin{equation}
 |g,n\rangle
\end{equation}
с системой в основном состоянии и два вектора
\begin{equation}
 |e_1,n-1\rangle,qquad |e_2,n-1\rangle
\end{equation}
с системой в поперечном дублете. Начальный диагональный блок имеет
собственные значения
\begin{equation}
 \frac{p_n}{3},\qquad\frac{p_{n-1}}{3},\qquad\frac{p_{n-1}}{3}.
\end{equation}
Энергосохраняющий унитарий вращает только внутри этого блока. По принципу
Ки Фана максимальный вес одномерного ground-подпространства равен
наибольшему собственному значению блока.

С учётом граничного сектора полной энергии ноль получаем точную оценку
\begin{equation}
 a_{\rm out}\leq
 \frac13\left[
 p_0+\max(p_1,p_0)+\max(p_2,p_1)+\max(p_3,p_2)
 \right].
 \label{eq:v6-four-clock-ground-bound}
\end{equation}
Правая часть является выпуклой кусочно-линейной функцией на симплексе.
Её максимум достигается на вершине и равен
\begin{equation}
 \boxed{a_{\rm out}\leq\frac23.}
 \label{eq:v6-four-clock-axis-ceiling}
\end{equation}

Машинный аудит проверил все четыре вершины и $10^5$ случайных состояний
часов. Максимум выборки равен $0.666061\ldots$ и не нарушает точную
границу.

\section{Кристаллическая фаза лежит выше потолка}

Канонический переход требует
\begin{equation}
 a_*=0.9121665963\ldots,
\end{equation}
поэтому дефицит минимальной энергосохраняющей модели равен
\begin{equation}
 a_*-\frac23=0.2454999296\ldots.
 \label{eq:v6-four-clock-axis-deficit}
\end{equation}
Значит, достаточная энтропийная ёмкость $\log4$ ещё не означает
достаточную энергетическую кратность.

Причина прозрачна. Две ортогональные поперечные компоненты должны отдать
одну и ту же энергию $\varepsilon_*$, не нарушая унитарности. Но на каждом
резонансном разрыве у невырожденных часов имеется только один вектор
$|g,n\rangle$, способный принять их амплитуду. Для полного переноса нужны
как минимум два ортогональных состояния часов на одном энергетическом
разрыве.

\begin{theorem}[Потолок невырожденных четырёхтактных часов]
Для одноосного гамильтониана \eqref{eq:v6-uniaxial-qutrit-hamiltonian} и
невырожденной четырёхуровневой лестницы любой унитарий, коммутирующий с
полной аддитивной энергией, переводит исходное $I_3/3$ в состояние с
выделенным собственным весом не более $2/3$. Требуемый вес равен
\begin{equation*}
 a_*=0.9121666\ldots,
\end{equation*}
поэтому каноническая одноосная фаза недостижима в минимальной модели.
\end{theorem}

\section{Почему четыре характера не дают нужного вырождения}

Проектный оператор порядка четыре имеет различные характеры
\begin{equation}
 1,\quad i,\quad-1,\quad-i.
\end{equation}
Следовательно, его четыре квазиэнергии различны по модулю периода. Выбор
других ветвей логарифма меняет абсолютные энергии на целые кванты
$2\pi/\tau$, но не создаёт двух состояний с одним и тем же характером.
Именно поэтому один четырёххарактерный носитель не поставляет требуемую
кратность два.

Разрушить вырождение поперечного дублета в $H_S$ нельзя без цены: тогда
орбита перестаёт быть точным $SO(3)/O(2)=\mathbb{RP}^2$ и механизм
проекторных дефектов теряет исходную симметрию.

\begin{nogocriterion}[Запрет скрытого расщепления дублета]
Нельзя обходить энергетический потолок назначением двух разных энергий
поперечным состояниям, одновременно продолжая считать вакуумное
многообразие точным $\mathbb{RP}^2$. Такое расщепление является новым
анизотропным входом.
\end{nogocriterion}

\section{Сопоставление с квантовой термодинамикой}

Тепловые операции по определению реализуются энергосохраняющими унитариями
с подходящими ресурсными состояниями. В конечных системах допустимость
перехода ограничивается не одной средней свободной энергией, а семейством
монотонностей; см. \path{doi:10.1038/ncomms3059},
\path{arXiv:1305.5278} и \path{doi:10.1073/pnas.1411728112}.
Когерентность между энергиями также является отдельным ресурсом и не
снимает автоматически популяционные ограничения; см.
\path{doi:10.1038/ncomms7383}.

Проектный результат является конкретной конечномерной версией этой
границы: общая энтропия разрешает переход, но симметрия энергии и
недостаточная кратность запрещают его минимальной машине.

\section{Строгие маршруты переоткрытия}

Гейт оставляет четыре, и только четыре, честных маршрута:
\begin{enumerate}
 \item найти в уже существующем $H_{15}$ или другом доказанном носителе
 две ортогональные копии одного резонансного характера;
 \item вывести второй резонансный канал часов;
 \item получить неаддитивную энергию взаимодействия из родительской
 суперсвязности;
 \item использовать специальные начальные корреляции, отдельно доказав
 их происхождение и ресурсный баланс.
\end{enumerate}

Первый маршрут минимален, поскольку не вводит новых состояний. Следующий
проверка должна провести аудит всех уже существующих кратностей $15$, $18$,
$20$, $35$ и четырёх характерных проекторов и проверить, является ли хотя
бы одна из них настоящим энергетическим вырождением, а не только
коэффициентной или семейной кратностью.

\begin{protocol}
\begin{enumerate}
 \item одноосный гамильтониан триплета:
 \textbf{$\varepsilon\operatorname{diag}(0,1,1)$};
 \item требуемый разрыв при сосуществовании:
 \textbf{$1.9664146236\ldots$};
 \item невырожденная четырёхтактная лестница:
 \textbf{проверена в резонансе};
 \item точный энергетический потолок выделенного веса:
 \textbf{$2/3$};
 \item требуемый выделенный вес: \textbf{$0.9121665963\ldots$};
 \item минимальные часы достигают фазы: \textbf{нет};
 \item причина: \textbf{нехватка резонансной кратности два};
 \item расщепление поперечного дублета: \textbf{запрещено симметрией};
 \item энтропийная ёмкость часов: \textbf{достаточна};
 \item энергосохраняющая структура часов: \textbf{недостаточна};
 \item следующая проверка:
 \textbf{аудит существующих кратностей как резонансного стока};
 \item рождение материи:
 \textbf{минимальным четырёхтактным холодильником не замкнуто}.
\end{enumerate}
\end{protocol}

Машинный сертификат сохранён в
\path{s2t_v6_clock_controlled_energy_conserving_quench_gate_results.json}.